\documentclass[book.tex]{subfiles}
\begin{document}

\chapter{Universal Differential Transformer}

\label{chap:deep_onet}
\noindent\rule{11cm}{0.4pt} \\
\textbf{Prerequisites:} \autoref{chap:dft} \\
\textbf{Difficulty Level:} ** \\
\noindent\rule{11cm}{0.4pt}
\vspace{5mm}

Idea: Pretrain a transformer on diff-eq (ODE/PDE) trajectories from many different types of PDEs. Use a mesh representation of the PDE data to allow for data from different mesh resolutions. Also include a basic representation of the transformer equation as one of the core data inputs to the transformer. The pretrained model should then be capable of predicting forward in time multiple different types of PDEs. 

This system could be useful then for multiscale physics. The transformer can jointly model a system such as a QM/MM system (say a protein enzyme with a catalytic region), or possibly a weather system. The joint transformer could also dynamically decide which regions to drop down to a higher resolution PDE simulation in order to run at larger timescales.

This system could also train on MD trajectories. The challenge would be to handle very long trajectories for training and to enforce correctness constraints.

\end{document}